\begin{textsample}{POS Dim 1 – ap – Score 44.00 – ap/ap\_inf\_e\_ef\_10.txt}  \label{ex:f1_pos_172}
Aporte \textbf{Metodológico} As situações, os perfis e as faixas etárias da modalidade EJA partilham expectativas e experiências diferenciadas.
Isto requer \textbf{uma} proposta educativa \textbf{que} alcance as expectativas e atenda às necessidades de adolescentes, jovens, adultos e idosos, reconhecendo suas especificidades de vivência corporal, sexual, cultural e de identidade.
De acordo com Teixeira ( 2006 ) : > \textbf{Educar} jovens e adultos para a vida é \textbf{um} desafio.
Repensar quais são os objetivos, as metas, os enfoques, as epistemologias, as \textbf{teorias} \textbf{que} fundamentam a \textbf{docência} não é \textbf{uma} \textbf{tarefa} fácil, mas necessária.
Precisa -se transformar a educação para transformar a realidade recursivamente, tornando a recíproca verdadeira. ( TEIXEIRA, 2006, p. 192 ).
Desta forma, esta Diretriz recomenda a construção coletiva ( comunidade escolar e outros ) de \textbf{uma} proposta pedagógica em cada unidade escolar, \textbf{que} contemple as especificidades e os segmentos da adolescência de 15 à 17 anos ; jovens de 18 à 29 anos ; adultos de 30 à 59 anos ; idosos à partir de 60 anos, como é referendado no Parecer das Diretrizes Curriculares para a EJA ( CNE/CEB nº 11/2000 ), \textbf{que} recomenda \textbf{atentar} -se para o perfil dos estudantes, suas experiências e situações reais de existência, devendo constituir a ideia \textbf{central} da organização desta modalidade ( BRASIL, 2000a ).
No contexto discutido na VI CONFITEA ( 2006 ), é importante também romper com a dicotomia estabelecida entre a formação para o trabalho e a formação de caráter geral.
Ambas merecem reflexão, pois olhar a EJA com \textbf{uma} educação somente voltada para o trabalho é reduzir ao tecnicismo da mão de obra, visão já superada \textbf{historicamente}.
Não se pode reduzir às necessidades do mercado de trabalho, nem tão pouco ser alheio às indigências de sobrevivência e às exigências da produção econômica, de onde os \textbf{sujeitos} sociais retiram sua sobrevivência, ou seja, a educação não deve ser refém da empregabilidade, nem de costas para o trabalho, deve ser recíproca ao mundo do trabalho ( mas \textbf{que} não seja ornamental, livresca ou oligárquica ) e intrínseca ao princípio de educação ao longo da vida.
Nas palavras de Ireland \textbf{et} al. ( 2005 ) : > O lugar do trabalho na vida do jovem e do adulto precisa ser o lugar do ser, onde ele se realiza \textbf{enquanto} produtor de cultura.
Articular -se ao mundo do trabalho não significa render -se ao utilitarismo e à empregabilidade, mas ter como princípio \textbf{central} a formação do trabalhador, na perspectiva do “ ser ” ( IRELAND \textbf{et} al., p. 38, 2005 ).
Os educandos da EJA possuem experiências sociais com o mundo da escrita, porque convivem com grupos diferenciados, seja na vizinhança, na igreja, na escola, na relação com os filhos e trabalho, alfabetizados ou não, mas possuem conhecimentos e funcionalidades da escrita, pois convivem em \textbf{uma} sociedade letrada.
Nestes termos e compartilhando as ideias de Freire ( 2003 ), é necessário \textbf{que} : > Antes de qualquer tentativa de discussão, de técnica, de materiais, de \textbf{métodos} para \textbf{uma} aula dinâmica assim, é \textbf{preciso}, indispensável mesmo, \textbf{que} o professor se ache “ repousado ” no saber de \textbf{que} a pedra fundamental é a curiosidade do ser humano.
É ela \textbf{que} me faz perguntar, conhecer, atuar, mais perguntar, reconhecer ( \textbf{FREIRE}, 2003, p. 52 ).
O melhor caminho é ouvir o \textbf{que} os educandos anseiam, suas dúvidas, interesses, \textbf{vontades}, curiosidades, porém sem distanciar e sem \textbf{perder} o foco das aprendizagens.
A EJA deve favorecer a aquisição de conhecimentos, habilidades, capacidades para a formação autônoma em prol de cidadãos aptos e ativos na sociedade.
Desta forma, o ensino da EJA deve reconhecer as características e especificidades de seus \textbf{sujeitos} \textbf{que} perpassam por aspectos como : as exigências e interesses da diversidade ; da infraestrutura \textbf{que} acolha a realidade do público ; a flexibilidade de tempo e espaços ; a disponibilidade de recursos didáticos \textbf{que} atendam e desenvolvam as potencialidades do \textbf{educando}.
O ambiente de aprendizagem da EJA \textbf{que} promova efetiva aprendizagem fundamenta -se, de acordo com Teixeira ( 2006 ), em cinco princípios : 1.
A aprendizagem deve ser centralizada em problemas e nas experiências ; 2.
As experiências devem ser significativas para o estudante ; 3.
O \textbf{aprendiz} deve ter liberdade de analisar a experiência ; 4.
As metas e as pesquisas devem ser fixadas e executadas pelo aluno, o estudante deve sentir -se livre de errar, de explorar alternativas para solução dos problemas e de participar nas decisões sobre a organização do seu ambiente de aprendizagem ; 5.
O aluno deve receber o"feed-back"sobre o seu progresso em relação às metas ( TEIXEIRA, 2006, p. 34 ).
A educação centrada no aprendedor, tendo como subsídio a Andragogia na definição de Malcolm Knowles, como ciência \textbf{que} orienta o adulto a aprender, pois não são \textbf{aprendizes} sem experiência, mas alunos da escola da vida, com o conhecimento \textbf{que} vem da realidade.
Desta forma, a andragogia tem como princípios : 1.
Autonomia : auto administrar -se e gosta de ser percebido e tratado como tal pelos outros. 2.
Experiência : oferece \textbf{uma} excelente base para o aprendizado de novos conceitos e novas habilidades. 3.
Prontidão para a Aprendizagem : maior interesse em aprender aquilo \textbf{que} está relacionado com situações reais de sua vida. 4.
Aplicação da Aprendizagem : as visões de futuro e tempo do adulto levam -no a favorecer a aprendizagem daquilo \textbf{que} possa ter aplicação \textbf{imediata}, o \textbf{que} tem como corolário \textbf{uma} preferência pela aprendizagem centrada em problemas em detrimento de \textbf{uma} aprendizagem centrada em áreas de conhecimento. 5.
Motivação para Aprender : são mais afetados pelas motivações internas \textbf{que} pelas motivações externas ( TEIXEIRA, 2006, p. 36 ).
Possuem experiências \textbf{que} passam durante a vida e o conhecimento \textbf{que} vem da realidade, buscam desafios e soluções \textbf{que} façam diferenças em suas vidas.
Eles aprendem melhor quando o assunto faz relação com a sua vida diária, diferencia -se dos demais, na consciência de \textbf{que} precisa do conhecimento, \textbf{que} este lhe faz falta ( WALL e TELLES, 2004 ).
Na Educação de Jovens e Adultos, a metodologia de ensino e aprendizagem fundamenta -se nos eixos cultura, trabalho e tempo, \textbf{que} são articuladores de toda ação pedagógico-curricular e da motivação e experiência dos \textbf{aprendizes}.
Os educandos trazem para o espaço-tempo escolar tanto a \textbf{marca} da destituição de direitos, quanto à riqueza de suas experiências de luta pela vida.
A concepção proposta para as ações pedagógicas desenvolvidas na EJA são \textbf{dialógicas}, reflexivas e críticas, voltadas para questões sociais, culturais, políticas, dentre outras, compreendendo o \textbf{educando} como sujeito \textbf{que} participa e interfere na construção histórica da sociedade em \textbf{que} \textbf{vive}.
Nesta perspectiva, Pinsky ( 2002 ) evidencia \textbf{que} : > O professor precisa conhecer as bases da cultura, o desenvolvimento do capitalismo, os movimentos sociais, as condições de vida das populações no passado, sua cultura material e suas ideias, a música de Beethoven, o cinema de Charles Chaplin, a literatura de Machado de Assis e por aí afora.
Noutras palavras, cada professor precisa, necessariamente, ter \textbf{um} conhecimento \textbf{sólido} do patrimônio cultural da \textbf{humanidade} ( PINSKY, 2002, p. 22 ).
As práticas pedagógicas devem privilegiar estratégias \textbf{que} contemplem as diferentes linguagens verbal ou alfabética e não verbal iconográfica ( leitura de imagens, desenhos, filmes, outdoors ) e cinética ( sonora, olfativa, tátil, visual e gustativa ), para \textbf{que} o \textbf{educando} reconheça as diferentes formas de falar, escrever e interpretar, bem como os efeitos dessas linguagens no processo comunicativo.
Em se tratando da Educação Socioeducativa e Penitenciária, devem garantir as seguintes competências : I – Pessoal : aprender a ser, ou seja, relaciona -se com a capacidade de conhecer a si mesmo, compreender -se, aceitar -se ; II – Social : aprender a conviver é a capacidade de relacionar -se de forma harmoniosa e produtiva com outras pessoas ; III – Produtiva : aprender a fazer, aquisição de habilidades necessárias para se produzir bens e serviços ; IV – Cognitiva : aprender a ser e a conviver, adquirir os conhecimentos necessários ao seu crescimento pessoal, social e profissional, assegurando a empregabilidade e/ou a trabalhabilidade a serviço da formação \textbf{baseada} em valores ( DELORS, 2012 ).
Desta forma, ao se pensar no processo educativo no espaço socioeducativo e penitenciário, deve -se ter a clareza sobre os limites ( tempo, espaço, cultura, entre outros ) impostos e as peculiaridades de seus \textbf{sujeitos} e locais.
Levando -se em consideração as especificidades da educação em espaços de privação de liberdade, deverão incentivar a promoção de novas estratégias pedagógicas, produção de materiais didáticos e a implementação de novas metodologias e tecnologias educacionais, assim como de programas educativos na modalidade Educação à Distância ( EAD ) ( BRASIL, 2010a ).
As Diretrizes Nacionais para a oferta de educação nos estabelecimentos penais orientam \textbf{que} a relação educação/trabalho, qualificação técnica e profissional do \textbf{educando} privado de liberdade, deve acontecer ainda durante o cumprimento da medida ou da pena e a integração deste à proposta de reabilitação com o trabalho, dispondo das qualificações e formações profissionais ofertadas, para uso dentro do próprio sistema de privação, como são os casos do monitor de educação e do agente prisional de saúde ( BRASIL, 2009a ).
Recomenda -se de forma geral, independente do público da EJA a ser beneficiado, o uso da \textbf{problematização} dos temas, pois possibilita ao \textbf{educando} a aprendizagem de novos conhecimentos, por meio da articulação entre os saberes e experiências acumulados e os saberes científicos, utilizando as várias linguagens textuais e as situações problemas como mediadoras do processo de construção individual e coletiva de conhecimento.
A pesquisa orientada também favorece possibilidades de compreensão da temática ; o diálogo entre educador e \textbf{educando} vão \textbf{revelar} o conhecimento prévio \textbf{que} possuem sobre o assunto.
Nesse contexto, os novos caminhos propõem a \textbf{interdisciplinaridade} como \textbf{uma} estratégia para compreensão, interpretação e explicação das temáticas, ou seja, o diálogo entre diversos saberes para ampliar a compreensão do objeto de estudo.
Parafraseando Minayo ( 2010 ), o produto final dessa estratégia será a \textbf{transdisciplinaridade}, “ \textbf{um} \textbf{método} diferente de \textbf{ver} as coisas e a vida ”.
Porém, há de se ressaltar \textbf{que} o acúmulo de saberes deve \textbf{despertar} e sensibilizar sobre consciência e a harmonia das relações entre seres \textbf{humanos} e natureza, \textbf{implicando} nas problemáticas ambientais \textbf{desencadeadas} pelo mau uso dos recursos naturais e tecnológicos, o \textbf{que} compromete a sustentabilidade socioambiental. \textbf{Demandam} -se assim, práticas \textbf{que} respeitem os espaços públicos, como bem coletivos da diversidade da vida e das culturas usados de forma democrática para sobrevivência das gerações do presente e futuras.
A participação e a ajuda mútua busca a construção da autonomia e da cooperação, cultivando valores como a solidariedade e o combate ao preconceito entre as formas de vida \textbf{que} coabitam o planeta.
Sugere -se a \textbf{contextualização} de temáticas integradas ao currículo e às atividades de Empreendimentos Econômicos Solidários ( organizados como cooperativas, associações, redes e outras formas ), nos quais os educandos/trabalhadores são os donos dos meios de produção e tomam decisões seguindo os princípios da autogestão.
A escola como espaço formador poderá colaborar na construção de \textbf{uma} sociedade sustentável, cabendo aos educadores, de acordo com o contexto da comunidade escolar nos quais estão inseridos, favorecer a formação de cidadãos críticos e atuantes, cientes da sua responsabilidade com a preservação do mundo.
O discurso do senso comum, histórico, político, econômico, filosófico, científico, entre outros, possibilitam a compreensão dos diversos pontos de vista para \textbf{que} o \textbf{educando} possa posicionar -se diante das diferentes situações do cotidiano.
A reflexão sobre as temáticas permite a \textbf{abordagem} dos conteúdos e articulação entre os eixos de cada \textbf{disciplina}, garantindo dessa forma a \textbf{interdisciplinaridade}.
Nesse sentido, a prática pedagógica deve facilitar a integração entre os diferentes saberes.
Espera -se \textbf{que} os educadores se identifiquem com este Documento e elaborem estratégias diferenciadas \textbf{que} sejam capazes de motivar, transformar o conhecimento e desenvolver habilidades⁸ e apresentem características peculiares tais como : ⁸ De investigação, de saber observar, de formular hipóteses, de buscar comprovações, de disposição à autocorreção, de \textbf{raciocínio}, formação de conceitos, de tradução ( RODRIGUES, 2010 ). - adolescente, jovem, adulto e idoso e as formas de inserção no mundo do trabalho ; - \textbf{espírito} de coletividade com vistas ao desenvolvimento de \textbf{um} trabalho pedagógico ; - visão global do currículo, postura inter/transdisciplinar e contextualizada, favorecendo o planejamento coletivo de estratégias pedagógicas ; - \textbf{percepção} do \textbf{educando} e de si mesmo como adultos em processo contínuo de formação ; - postura investigativa na prática educativa ; - \textbf{compromisso} ético e político com a \textbf{dignidade} humana ( FÓRUM EJA, 2005 ).

% matched lemmas: abordagem, aprendiz, atentar, basear, central, compromisso, contextualização, demandar, desencadear, despertar, dialógico, dignidade, disciplina, docência, educar, enquanto, espírito, et, freire, historicamente, humanidade, humanos, imediato, implicar, interdisciplinaridade, marca, metodológico, método, percepção, perder, preciso, problematização, que, raciocínio, revelar, sujeito, sólido, tarefa, teoria, transdisciplinaridade, um, ver, viver, vontade
\end{textsample}
